\documentclass{article}
\usepackage[dblblindworkshop, final]{neurips_2026}

\usepackage[utf8]{inputenc}
\usepackage[T1]{fontenc}
\usepackage{hyperref}
\usepackage{url}
\usepackage{booktabs}
\usepackage{amsfonts}
\usepackage{microtype}
\usepackage{xcolor}

\hypersetup{colorlinks=true,linkcolor=black,citecolor=black,urlcolor=black}

\workshoptitle{MATH-AI: The 6th Workshop on Mathematical Reasoning and AI}

\author{%
  Anirudh Balaji \\
  PES University, Bengaluru, India \\
  \texttt{PES1UG25AM042@stu.pes.edu} \\
}

\title{Structural Coverage: Do AI Provers Use Mathlib\\the Way Mathematicians Do?}

\begin{document}
\maketitle

\begin{abstract}
A Lean proof is not only a binary success. It is a selection of premises from a
library, and that selection induces a subgraph of mathlib4's dependency graph.
Two proofs can both be correct and induce structurally different subgraphs;
\texttt{pass@k} is blind to the difference. We define structural coverage over
these subgraphs and compare 29{,}750 released machine Lean proofs against
70{,}086 human mathlib declarations. Three findings. (i)
Machine proofs draw on a small subset of the human premise vocabulary, one to two
orders of magnitude smaller: 78{,}642 against 1{,}127 premises, Jaccard $0.0129$,
leaving $98.7\%$ of the human vocabulary untouched. The gap
survives every control we build, narrowing to $19\times$--$28\times$ when the
human corpus is restricted to the regions the machine corpus occupies and to
$5.7\times$ under an exploratory genre-matched comparison against mathlib's own
competition archive. (ii) The machine
corpus concentrates on a stratum that is essentially absent from mathlib:
arithmetic-tactic proofs of at most two tactic steps are $23.1\%$ of the machine
corpus and $0.02\%$ of mathlib, and a single premise, \texttt{sq\_nonneg},
accounts for $27.0\%$ of all machine premise occurrences. (iii) A cautionary
negative result: the graph's published depth attribute proved to be a centrality
proxy, not a derivation-length measure, and we withdraw a finding built on it.
Structural coverage is a non-redundant evaluation axis, cheap to compute from
released corpora.
\end{abstract}

\section{Introduction}

The MATH-AI call asks: \emph{``Humans vs.\ machines: What are the comparative
strengths, limitations, and characteristic failure modes of human mathematicians
and AI systems?''} Formal theorem proving offers an unusually direct way to
answer it. When a prover closes a goal in Lean, it does not merely produce a
truth value. It selects premises from mathlib4, and that selection induces a
subgraph of the library's dependency graph.\footnote{All metrics come from
committed CSVs produced by scripts with pinned data revisions and fixed seeds.
Code, the committed CSVs, the pinned revisions and the seeds will be released.}
Two proofs of the same statement can both compile and still induce subgraphs
that share almost nothing. Standard
evaluation cannot see this. \texttt{pass@k} counts closures; it is by
construction blind to which part of the library was used to get there.

We ask whether machine-generated proofs use mathlib the way human contributors
do, and answer by measurement. Our object of study is the \emph{explicit
subgraph} of mathlib4's dependency graph \citep{li2026network} --- the
source-visible portion approximating human-intended dependencies --- and our
comparison is between the premise vocabularies human and machine proofs induce
over it. Three findings follow.

\textbf{1. Vocabulary coverage differs by one to two orders of magnitude, and
the gap is not an artifact of corpus size or library area.} Across full corpora
the human premise vocabulary is $78{,}642$ against the machine corpus's
$1{,}127$, a ratio of $69.8\times$, with Jaccard $0.0129$ and $98.7\%$ of the
human vocabulary never touched by any machine proof. Every control we build narrows the gap and
none closes it: rarefaction to equal $n$ leaves $44.0\times$, a 1:1 class- and
step-matched design $18.8\times$, and restriction to the $42$ library regions the
machine vocabulary inhabits $28.5\times$.

\textbf{2. There is a characteristic concentration, and it is a stratum.}
Proofs classified as arithmetic-tactic with at most two tactic steps are
$23.1\%$ of the machine corpus and $0.02\%$ of mathlib, or $14$ of $70{,}086$
declarations. Machine premise usage is correspondingly concentrated:
\texttt{sq\_nonneg} alone is $27.0\%$ of machine premise occurrences and the top
25 premises are $68.1\%$.

\textbf{3. A negative methodological result about the released graph itself.}
We built a depth metric on the graph release's published per-declaration layer
attribute, promoted the resulting gap to a stated counter-finding, and then
found the attribute to be a centrality proxy. We withdraw the finding and report the failure, because it is a
second independent instance of a caution the graph's own authors raise
\citep{li2026network} and because a source-visible pipeline can consume such a
column for weeks without any signal that it is wrong.

\section{Setup}

\paragraph{The graph and the explicit subgraph}
We use the published mathlib4 dependency graph of \citet{li2026network}
($308{,}129$ declarations, $8{,}436{,}366$ edges) at a pinned revision, built on
mathlib4 at commit \texttt{534cf0b}. We do not rebuild it. That work reports
that $74.2\%$ of dependency edges are compiler-synthesized and invisible in
source, and defines an \emph{explicit subgraph} as a closer proxy for
human-intended dependencies. The explicit subgraph, $2{,}178{,}534$ edges, is
our object of study, and source-visible extraction is the appropriate instrument
for it.

\paragraph{Corpora}
Our human reference is the $70{,}086$ mathlib declarations we can locate in
source and extract premises from. Our machine corpus is
\texttt{Goedel-LM/Lean-workbook-proofs}, $29{,}750$ released proofs of Lean
Workbook problems, pre-verified by its releasers, so no Lean execution is needed
for correctness. Both are censuses, not samples.

\paragraph{The confound we cannot remove}
The machine corpus targets competition statements while mathlib targets library
development. No control available to us removes this confound on problem
selection; Section~\ref{sec:lim} states what it permits us to claim. The
comparison is over the shared premise vocabulary.

\section{Method}\phantomsection\label{mth}

\paragraph{Extraction}
We resolve identifier tokens in a declaration's proof body against the graph's
declaration names at declaration scope, comments and strings stripped, by exact
name or by completion under a namespace prefix in scope; bare dot-notation
residue (\texttt{.mp}) is kept separate, as in \citet{cslibpremise}.

\paragraph{Retention, and why we report it rather than gate on it}
Against explicit-subgraph ground truth on the mathematical hub layer the
extractor recovers $0.6124$ of targets. Scope decomposes that figure. Matching premise
names anywhere in the whole source file recovers only $0.8193$ of
mathematical-layer explicit edges, so an absolute target near $1$ is unreachable
by construction: that ceiling belongs to Lean's elaborator, not to our parser.
Nor is it the bound the extractor works against. At declaration scope, where it
runs, the reachable bound is $0.6948$, so efficiency is
$0.6124/0.6948 = \mathbf{0.8815}$; dividing by the whole-file $0.8193$ instead
gives $0.7475$, against a scope the extractor never sees. We report retention
with this decomposition, having mis-specified a pass/fail threshold twice.

\paragraph{The hub-layer partition}
Results below split premises into language-infrastructure and
mathematical-content layers. This partition is ours and provenance-based: a
declaration is infrastructure if its \texttt{file\_module} lies outside Mathlib,
if it is an instance or coercion, if its \texttt{kind} is a kernel kind, or if
its name carries a generated-name marker. It is unrelated to the per-declaration
layer attribute withdrawn in Section~\ref{sec:lim}.

\paragraph{Population neutrality}
What matters is not how much the instrument misses but whether it misses
differentially across corpora. Per-class efficiency against the reachable bound
is flat across the five tactic classes ($0.874$--$0.896$), so the instrument does
not preferentially penalise either corpus's tactic mix. That flatness is what
carries the claim, and we rest it on nothing else.

\paragraph{Length control}
Short proofs mechanically use fewer, commoner premises, so comparisons are
length-controlled on \emph{tactic steps}, not characters: the corpora overlap
$0.856$ on steps against $0.215$ on characters, so a character axis is largely a
formatting control on hand-wrapped mathlib source. Coverage needs no such
standardization: the effect lies orders of magnitude outside the regime a length
correction operates in, so no reweighting available to us could produce or
remove it. Corpus size is handled by design instead, in the rarefied and matched
rows of Table~\ref{tab:coverage}.

\section{Results}\phantomsection\label{res}

\subsection{Vocabulary coverage}

Table~\ref{tab:coverage} gives the primary result under five designs of
increasing severity. Two one-sided extractor artifacts --- a declaration
resolving its own name, and the \texttt{lemma} keyword resolving to a real
mathlib declaration --- inflated the human side only. They account for
$33.6\%$ of the uncorrected human count, and are removed throughout, taking the headline from $118{,}514$ to $78{,}642$ human
premises before any comparison is made.

\begin{table}[t]
\centering
\small
\begin{tabular}{lrrrrr}
\toprule
design & human & machine & ratio & Jaccard & human-only \\
\midrule
full corpora & 78{,}642 & 1{,}127 & $69.8\times$ & 0.0129 & 0.9871 \\
rarefied to machine $n$ & 49{,}601 & 1{,}127 & $44.0\times$ & --- & --- \\
\textbf{area-restricted (primary control)} & \textbf{32{,}091} & \textbf{1{,}127} & $\mathbf{28.5\times}$ & \textbf{0.0315} & \textbf{0.9684} \\
1:1 class- and step-matched & 21{,}162 & 1{,}123 & $18.8\times$ & 0.0360 & 0.9634 \\
genre-matched (exploratory) & 1{,}101 & 193 & $5.7\times$ & --- & --- \\
\bottomrule
\end{tabular}
\caption{Premise vocabulary under five designs. The area-restricted row keeps
the 42 regions the machine vocabulary occupies ($40.8\%$ of human premises, no
machine loss). The genre-matched row is exploratory; see
Section~\ref{sec:lim}.}
\label{tab:coverage}
\end{table}

The corpus-size objection is answered by measurement. Rarefying the human corpus
to the machine corpus's proof count still gives $49{,}601$ premises
($44.0\times$ at equal $n$), and the 1:1 design matched on tactic class and step
count gives $18.8\times$ with $96.3\%$ untouched.

The area-restricted design is our primary control, its rule fixed before it ran:
every top-level region containing at least one machine premise. The rule is
support-based, so there is no threshold or top-$k$ to tune. The machine vocabulary occupies $42$ of the
$1{,}595$ regions the human vocabulary spans; those hold $40.8\%$ of human
premises; inside them the human vocabulary is still $28.5\times$ larger, $96.8\%$
untouched. A tighter twelve-namespace variant gives $19.2\times$ but discards $31.0\%$ of
the machine vocabulary, so we report the range $19\times$--$28\times$.

\subsection{A characteristic concentration}

The machine corpus concentrates on a stratum essentially absent from mathlib.
Proofs with an arithmetic-flavoured tactic set (\texttt{nlinarith},
\texttt{linarith}, \texttt{positivity}, \dots) and at most two tactic steps are
$23.1\%$ of the machine corpus ($6{,}869$ of $29{,}750$) and $0.02\%$ of mathlib,
or $14$ of $70{,}086$ declarations. We report this as a result: it is the largest
cell of the machine corpus and has no human counterpart. Usage is concentrated to
match. \texttt{sq\_nonneg} alone is $27.0\%$ of machine premise occurrences and
the top 25 are $68.1\%$, and machine proofs resolve $2.08$ premises per proof
against $6.24$ for human declarations. The composition of that vocabulary differs too: $25.6\%$ of the
machine premise vocabulary sits in the language-infrastructure hub layer against
$4.6\%$ of the human one, $5.6\times$ more infrastructure-weighted. That is
consistent with the single-step automation stratum above: proofs closed by a
single arithmetic tactic draw disproportionately on infrastructure rather than
mathematical content.

\subsection{Structural breadth}

Breadth counts distinct Louvain communities touched per tactic step. Its sign is
stable and negative: machine proofs touch fewer regions per step at every scope.
Its magnitude is design-dependent by $4\times$, at $-0.607$ full corpora,
$-0.780$ common support, $-0.201$ matched 1:1. We quote all three, never
one alone. The partition is also coarse here (median community size $1$; the five
largest hold $68\%$ of the graph), so breadth is a secondary characterization
with its limitation attached, not a load-bearing result.

\paragraph{What this buys an evaluation}
Coverage separates proof sets that \texttt{pass@k} cannot, needs no
standardization or imputation, and is computable from any released corpus
against a published graph. Its Jaccard holds at $0.0129$ on full corpora and
$0.0128$ under common support, the restriction that reverses or destroys every
other comparison here; the 1:1 matched design moves it to $0.0360$.

\section{Limitations}\phantomsection\label{lim0}
\label{sec:lim}

\paragraph{Genre, and the control we do not have}
Area restriction fixes library \emph{area} but not \emph{genre}: a mathlib
theorem is general-purpose infrastructure, a competition problem a self-contained
puzzle. The proper control is a same-problem comparison and it is \textbf{not
currently available}: the only Lean~4 human miniF2F~\citep{minif2f} proofs in existence are $67$,
and they are a prover's own few-shot prompt pool; the standard Lean~4 port's
$488$ statements are all \texttt{sorry}, and the original release's $155$ proofs
are Lean~3. As an exploratory substitute, mathlib's \texttt{Archive/Imo} ($508$ human
declarations, no \texttt{sorry}) draws $1{,}101$ premises against $193$ for the
machine corpus rarefied to $508$ proofs --- $5.7\times$. That figure is not
pre-registered and is a lower bound: $467$ of the $508$ cite local lemmas the
extractor cannot see, \texttt{Archive} lying outside the graph's build target,
which also leaves it not extraction-audited.

\paragraph{Instrument}
Retention is $0.6124$. We argue population neutrality from flat per-class
efficiency rather than assume it, and under-recovery is conservative in
direction, but the vocabularies are instrument-filtered and should be read so.
Three further choices run the same way: \texttt{symm}/\texttt{trans} are left
unexcluded though they resolve $5{,}222$ times on the human side and $14$ on the
machine one, prefix resolution is uncorrected though it supplies $3.657$
premises per proof to the human side against $0.374$ to the machine one, and
coverage is reported over both hub layers, where the mathematical layer alone
would give $89.6\times$ against the reported $69.8\times$. All three leave the
reported gap smaller than a stricter accounting would give.

\paragraph{Depth is withdrawn}
We built a depth metric on the release's per-declaration layer attribute, which
we took to be distance from primitives. It is not, and the error was ours: its
lowest layer holds $119{,}633$ nodes of maximum in-degree $0$, uncited leaves
rather than axioms, its highest holds \texttt{Eq.refl}, \texttt{Set} and
\texttt{LE.le}, and it correlates $+0.677$ with log in-degree. That is what the attribute
should look like if edges run declaration-to-premise with sources at level zero,
so it is consistent with the release's definition; we read it backwards. As we use it the quantity is a centrality proxy,
so a gap in it cannot carry a depth claim, which is the inference
\citet{li2026network} rule out. We withdraw rather than reinterpret it, since
re-pointing a quantity at a new claim after its comparison has been seen is a
researcher degree of freedom. The check is cheap: correlate a derived
attribute against in-degree and inspect its extremes.

\paragraph{Scope}
One machine corpus. These are comparisons of two distributions, not causal
claims: we write ``machine proofs concentrate on,'' never ``provers are biased
toward.''

\section{Related work}\phantomsection\label{rel}

\citet{li2026network} give the multilayer analysis of mathlib4 we build on,
release the graph, and pose structural drift from AI proofs as future work; we
extend their human baseline to the machine population. \citet{theorembench} find
provers biased toward easy subtheorems and \citet{proofrank} assess proof quality
in ProofRank, both at theorem or whole-proof level, where we work on the induced
premise subgraph. \citet{verisoftbench} correlate success with transitive
dependency closure size, which counts how much a proof depends on, not where in
the library it sits.\phantomsection\label{repro}

\phantomsection\label{bibstart}
\small
\bibliographystyle{plainnat}
\begin{thebibliography}{9}

\bibitem[Li et al.(2026)]{li2026network}
Xinze Li, Nanyun Peng, Simone Severini, and Patrick Shafto.
\newblock The Network Structure of Mathlib.
\newblock \emph{arXiv:2604.24797}, 2026.

\bibitem[Pham et al.(2026)]{theorembench}
QuocViet Pham, Elvir Karimov, Andrey Galichin, and Ivan Oseledets.
\newblock TheoremBench: Evaluating LLMs on Theorem Proving in Formal
Mathematics.
\newblock \emph{arXiv:2606.09450}, 2026.

\bibitem[Petrov et al.(2026)]{proofrank}
Ivo Petrov, Jasper Dekoninck, Dimitar I. Dimitrov, and Martin Vechev.
\newblock Not All Proofs Are Equal: Evaluating LLM Proof Quality Beyond
Correctness.
\newblock \emph{arXiv:2605.10379}, 2026. ICML 2026 AI for Math workshop.

\bibitem[Xin et al.(2026)]{verisoftbench}
Yutong Xin, Qiaochu Chen, Greg Durrett, and I\c{s}il Dillig.
\newblock VeriSoftBench: Repository-Scale Formal Verification Benchmarks for
Lean.
\newblock \emph{arXiv:2602.18307}, 2026.

\bibitem[Ji(2026)]{cslibpremise}
Junye Ji.
\newblock CSLibPremiseBench: Structure-Guided Premise Retrieval and Label
Robustness for Lean 4 Computer-Science Theorems.
\newblock \emph{arXiv:2605.14549}, 2026.

\bibitem[Zheng et al.(2022)]{minif2f}
Kunhao Zheng, Jesse Michael Han, and Stanislas Polu.
\newblock MiniF2F: A Cross-System Benchmark for Formal Olympiad-Level
Mathematics.
\newblock \emph{ICLR}, 2022.

\end{thebibliography}

\end{document}
